\documentclass[a4paper,10pt,twoside,openany]{article}

\usepackage[lang=hebrew]{maths}
\usepackage{hebrewdoc}
\usepackage{stylish}
\usepackage{lipsum}
\let\bs\blacksquare

\usepackage{siunitx,array}
\sisetup{table-format=-1.5, group-digits=false}
\newcolumntype{L}{>{\phantom{$-$}}l}       % prefix some whitespace
\newcommand\mcL[1]{\multicolumn{1}{L}{#1}} % handy shortcut macro

\setlength{\parindent}{0pt}

%%%%%%%%%%%%
% Styling %
%%%%%%%%%%%%

\usepackage{enumitem}

%%%%%%%%%%%%%
% Counters  %
%%%%%%%%%%%%%

\setcounter{section}{1}     
            
%BIBLIOGRAPHY
\usepackage[
backend=biber,
style=alphabetic,
]{biblatex}
\addbibresource{bibliography.bib} %Imports bibliography file

%%%%%%%%%%
% Title  %
%%%%%%%%%%
\title{
אלגברה ב' - גיליון תרגילי בית 8 \\
צורת ז'ורדן
\\
\vspace{1cm}
\large{תאריך הגשה: 25.6.2026}
}
\date{}

\begin{document}
\maketitle

\begin{exercise}
עבור ההעתקות הבאות $f_i$, קיבעו האם $f_i$ מכפלה פנימית. הוכיחו את קביעתכן.

\begin{enumerate}
\item
\begin{align*}
f_3 \colon \Mat_n\prs{\mbb{R}} \times \Mat_n\prs{\mbb{R}} &\to \mbb{R} \\
\prs{A,B} &\mapsto \tr\prs{B^t A}
\end{align*}

\item
\begin{align*}
f_3 \colon \Mat_n\prs{\mbb{C}} \times \Mat_n\prs{\mbb{C}} &\to \mbb{C} \\
\prs{A,B} &\mapsto \tr\prs{B^t A}
\end{align*}

\item
\begin{align*}
f_4 \colon \mbb{R}_{\leq n}\brs{x} \times \mbb{R}_{\leq n}\brs{x} &\to \mbb{R} \\
\prs{f,g} &\mapsto f\prs{0} g\prs{0} + \ldots + f\prs{n} g\prs{n}
\end{align*}

\item
\begin{align*}
f_5 \colon \mbb{C}_{\leq n}\brs{x} \times \mbb{C}_{\leq n}\brs{x} &\to \mbb{C} \\
\prs{f,g} &\mapsto f\prs{0} g\prs{0} + \ldots + f\prs{n} g\prs{n}
\end{align*}
\end{enumerate}
\end{exercise}

\begin{exercise}
היעזרו באי־שוויון קושי־שוורץ כדי להראות שמתקיים
\[\text{.} \forall x,y,z \in \mbb{R}_+ \colon x + y + z \leq 2 \prs{\frac{x^2}{y+z} + \frac{y^2}{x+z} + \frac{z^2}{x+y}}\]
\textbf{רמז:}
חישבו כיצד לפרש את אגף ימין בעזרת נורמה.
\end{exercise}

\begin{exercise}
יהי
$V = \mbb{R}_2\brs{x}$
ותהיינה
\begin{align*}
\trs{f,g}_1 &= \int_0^1 f\prs{x} g\prs{x} \diff x \\
\trs{f,g}_2 &= f\prs{-1} g\prs{-1} + f\prs{0} g\prs{0} + f\prs{1} g\prs{1}
\end{align*}
שתי מכפלות פנימיות על
$V$.
יהי
\[\text{.} W = \set{f \in V}{f\prs{x} = f\prs{-x}} \leq V\]
\begin{enumerate}
\item 
מיצאו בסיס
$B = \prs{w_1, \ldots, w_m}$
של
$W$
והשלימו אותו לבסיס
$C$
של
$V$.
בצעו את תהליך גרם־שמידט על
$C$
לפי כל אחת מהמכפלות הפנימיות כדי לקבל בסיסים אורתונורמליים לפיהן.

\item
היעזרו בבסיסים שמצאתן בסעיף הקודם כדי למצוא
את
\[W^\perp \coloneqq \set{v \in V}{\forall w \in W: \trs{v,w} = 0}\]
לפי כל אחת מהמכפלות הפנימיות.

\item
מיצאו את ההטלה
$P_W$
על
$W$
במקביל ל־%
$W^\perp$
לפי כל אחת מהמכפלות הפנימיות.
הטלה זאת נקראת
\emph{ההטלה האורתוגונלית על
$W$}
(לפי המכפלה הפנימית המתאימה) ונראה בהמשך כי
$P_W\prs{v}$
הוא הוקטור הקרוב ביותר ל־%
$v$
ב־%
$W$.
\end{enumerate}
\end{exercise}

\end{document}